\section{附录1：程序代码}

本文所有程序代码环境为 Python 3.13 。代码文件清单如表\ref{tab:code_files}所示：

\begin{table}[htbp]
    \centering
    \caption{程序文件清单}
    \label{tab:code_files}
    \small
    \begin{tabular}{clp{7cm}}
        \toprule
        序号 & 文件名 & 用途 \\
        \midrule
        1 & \texttt{style\_config.py}   & 全局绘图风格与统一配色方案 \\
        2 & \texttt{process.py}         & 数据清洗与探索性数据分析 \\
        3 & \texttt{model.py}           & 特征工程、三模型训练、评估与交叉验证 \\
        4 & \texttt{shap\_analysis.py}  & SHAP 归因分析与可视化 \\
        5 & \texttt{retention.py}       & 客户风险分层与挽留策略 ROI 计算 \\
        \bottomrule
    \end{tabular}
\end{table}

\subsection*{全局风格配置}
\lstinputlisting[
    language=Python,
    caption={style\_config.py —— 全局绘图风格与配色},
    label=lst:style,
    basicstyle=\scriptsize\ttfamily,
]{code/style_config.py}

\subsection*{探索性数据分析}
\lstinputlisting[
    language=Python,
    caption={process.py —— 数据清洗与 EDA 可视化},
    label=lst:process,
    basicstyle=\scriptsize\ttfamily,
]{code/process.py}

\subsection*{逻辑回归建模与模型验证}
\lstinputlisting[
    language=Python,
    caption={model.py —— 特征工程、三模型训练、交叉验证与阈值分析},
    label=lst:model,
    basicstyle=\scriptsize\ttfamily,
]{code/model.py}

\subsection*{SHAP 可解释性分析}
\lstinputlisting[
    language=Python,
    caption={shap\_analysis.py —— SHAP 归因分析与可视化},
    label=lst:shap,
    basicstyle=\scriptsize\ttfamily,
]{code/shap_analysis.py}

\subsection*{客户风险分层与 ROI 计算}
\lstinputlisting[
    language=Python,
    caption={retention.py —— 分位数分层与挽留策略 ROI 估算},
    label=lst:retention,
    basicstyle=\scriptsize\ttfamily,
]{code/retention.py}
